\section{Introduzione dell'utilizzo del cluster THOR}
\label{sec:thor}

Il 10 luglio 2026 il lavoro riprende dopo una decina di giorni di pausa (l'ultimo commit precedente, \texttt{1b6a7e2}, risale al 30 giugno) con \texttt{4bf9253} (\emph{``RESTORED: previous gpt working version''}), che aggiunge un nuovo file \texttt{2\_1\_benchmark\_gpt.ipynb} di 415 righe, di contenuto identico byte-per-byte al notebook \texttt{2\_benchmark.ipynb} del commit iniziale \texttt{0bf3213} --- e non alla versione successivamente modificata in \texttt{1b6a7e2}. Il repository non permette di stabilire con certezza il motivo esatto di questo ripristino (ad es. se una modifica locale non fosse più recuperabile, o se si trattasse di una scelta deliberata di ripartire da una base nota); è però verificabile che si tratti di un ripristino puntuale della versione più antica salvata, non di una modifica incrementale.

Pochi minuti dopo, \texttt{a366462} (\emph{``ADDED: first configuration of Thor Server Files''}) rinomina il notebook in \texttt{2\_1\_benchmark\_thor.ipynb} e lo modifica per puntare a un server LLM locale (\texttt{base\_url="http://localhost:8080/v1"}) che serve il modello \texttt{qwen2.5-7b-instruct} tramite \texttt{llama.cpp}, sostituendo la chiamata cloud a OpenAI usata fino a questo punto. Il giorno successivo, \texttt{52cd3d8} (\emph{``ADDED: new Thor setup files added''}) aggiunge l'intera infrastruttura di supporto in \texttt{setup\_thor/}: le definizioni Apptainer per i due container (\texttt{llama\_server.def}, \texttt{jupyter\_env.def}), lo script di lancio del server (\texttt{llama\_server.sbatch}), un primo script di esecuzione automatica del notebook (\texttt{run\_benchmark.sbatch}), una patch per la connessione dinamica al server (\texttt{notebook\_patch.py}) e due guide operative (\texttt{README.md}, \texttt{thor-getting-started.md}, quest'ultima di oltre 600 righe) più una guida alternativa basata su vLLM (\texttt{guida\_vllm\_thor.md}).

La configurazione messa in piedi in questi due commit permette un primo esperimento funzionante su Thor (analizzato in \S\ref{sec:limiti}), ma --- come emerso da una sessione di debug sistematica il 20 luglio 2026, un mese e mezzo dopo --- restava di fatto poco collaudata end-to-end: nomi di job non corrispondenti tra script e job realmente in esecuzione, richieste di GPU su una partition per cui l'utente non aveva permessi, un problema di risoluzione utente (NSS/LDAP) del cluster specifico dei nodi \texttt{compute} che faceva fallire silenziosamente sia \texttt{apptainer} sia i comandi Slurm all'interno di job batch, e dipendenze Python installate senza vincolo di versione che si sono rivelate incompatibili con il codice del progetto. Questa sessione di debug (i commit da \texttt{40e1fbd} a \texttt{549d104} e il file \texttt{setup\_thor/COME\_LANCIARE\_NOTEBOOK.txt} ne sono la registrazione diretta) ha prodotto la prima esecuzione automatica riuscita dell'intera pipeline end-to-end su Thor.
